% !PRINTED 608
\weberpage{608}
Nun ist nach \S.\,164, (8) die Differenz $\Theta - \Theta'$ gleich der Quadratwurzel aus der Grundzahl von $\Omega$, so dass wir aus (6) und (8) erhalten:
\weq{9}{\omega = \frac{b+\sqrt{\Delta}}{2c}, \quad \Delta = b^2 + 4ac,}
worin das Vorzeichen von $\sqrt{\Delta}$ durch die in \S.\,164, (6) gemachte Annahme \"uber das Vorzeichen von $\sqrt{d}$ bestimmt ist.

Es ist also $\omega$ eine zur Discriminante $\Delta$ geh\"orige \gesperrt{quadratische Irrationalzahl}, wie wir sie im \S.\,122 des ersten Bandes betrachtet haben, und die Form $\varphi$ kann, wenn wir $a_{2,2}=e$ setzen, so ausgedr\"uckt werden:
\weq{10}{\varphi = ec(t_1+\omega t_2).}

Hierin ist $c$ die kleinste positive ganze rationale Zahl, f\"ur die das Product $c\omega$ eine ganze Zahl wird, und jede andere ganze rationale Zahl $k$, durch die das Product $k\omega$ eine ganze Zahl wird, muss durch $c$ theilbar sein.

Man erkennt dies am einfachsten, wenn man $\omega$ nach \S.\,164, je nachdem $b$ gerade oder ungerade ist, in eine der beiden Formen setzt:
\weq{11}{\omega = \frac{1}{c}\left(\frac{b}{2}+\Theta\right), \quad \frac{1}{c}\left(\frac{b-1}{2}+\Theta\right),}
woraus, da $1, \Theta$ eine Basis von $\mathfrak{o}$ ist, das Gesagte hervorgeht.

Darauf gr\"undet sich nun der Beweis, dass, wenn $\omega$ irgend eine zur Discriminante $\Delta$ geh\"orige quadratische Irrationalzahl bedeutet, $c$ die kleinste positive ganze rationale Zahl, die $c\omega$ zu einer ganzen Zahl macht, und $e$ eine beliebige positive ganze rationale Zahl, auch umgekehrt durch
\weq{12}{\varphi = ec(t_1+\omega t_2)}
immer eine Basisform dargestellt ist.

Dazu haben wir nur zu zeigen, dass
\[
\alpha_1 = ec, \quad \alpha_2 = ec\omega
\]
eine Basis des durch (12) definirten Functionals $\varphi$ ist.

Bilden wir zu diesem Zwecke zun\"achst die Norm von $\varphi$, so erhalten wir nach (7)
\weq{13}{N(\varphi) = e^2 c(ct_1^2 + bt_1t_2 - at_2^2),}
und folglich die absolute Norm
\weq{14}{N_a(\varphi) = e^2 c.}
